\section{When do the one- and two-letter lifts suffice?}

% BEGIN H83F1
\paragraph{H83F1. Full local and global bounds in the absence of adjacent thick positions.}
At a mixed inert C or D place, assume its cyclic word has no adjacent
pair of thick positions. Then the required bound holds on every
tested finite-cohomology map. This assertion has no upper or lower
prime-size restriction beyond the original guard. Indeed, cutting
immediately before each nonthick position gives arcs of length one
or two. The length-one scalar lifts and the length-two lifts ending
in an up letter are in the signed scalar chapter. For an arc ending
in a down letter, that chapter treats the large initial filtered
index and D83L2 treats its complement, including the band. Each
result is within the common integer allowance
$E_1=\lfloor40U+\eta I\rfloor$, since the exponents are integers,
$U\ge1$, $R,\sigma\le3U$, and the length-two bound is
$15U+1+\eta I$. A lift at a smaller integer exponent can be
multiplied by the missing power of $p$. S83M2 now assembles those
lifts with the at most $4dm$ exceptional generators, using the
actual target classes of U83M6. The resulting normalized length
fits $\mathrm{lenRowCD}$ by U83M4.

At any given index, a completely nonthick word is covered by
N82I1, and an entirely thick word by T83A6. Thus if no mixed
orbit at that index has a double thick hit, use G82I2 to obtain
the native global pointwise inequality for every good point. In
particular this completes every guarded inert row D: when the
second-component order is small the word criterion in I82W5
applies, and in full generality M82R2 covers row D at
$\psi\ne\mathrm{cc}$ while T83A6 covers $\psi=\mathrm{cc}$.
% END H83F1

% BEGIN H83F2
\paragraph{H83F2. An explicit finite test on the Frobenius increment.}
Let $v$ be the additive second component of inverse Frobenius,
$a_0$ the element of order $n_0$ in I82W4, and $T_2$ the
member's actual host digit set. Define the ordered-pair count
\[
   M(v)=\#\{(s,t)\in T_2^2:s\ne t,\ (t-s)a_0=v\}.
\]
At every guarded middle index with $\psi\ne\mathrm{cc}$ and
$M(v)=0$, the native inequality holds for every good point,
at all eligible prime sizes. The same holds whenever
$\operatorname{ord}(v)\le4 L_w$.

Indeed the second coordinates of two consecutive thick letters are
$u_0+s a_0$ and $u_0+t a_0$. Equal digits specify the same
underlying real embedding. Consecutive embeddings above that real
embedding differ by $1$ or complex conjugation. Inertia-free
Frobenius here has negative quadratic sign, and the latter element
was excluded. Thus the digits would be distinct and contribute
to $M(v)$. If $k=\operatorname{ord}(v)\le4 L_w$, such a pair
would give $n_0\mid k(t-s)$, although $0<k|t-s|<n_0$ by the
strict digit bounds. In both cases H83F1 applies.

In particular the only second components for which this test can require
the length-three argument lie among the nonzero differences
$(t-s)a_0$, with
\[
 0<|t-s|\le\left\lfloor\frac{n_0-1}{4L_w}\right\rfloor .
\]
Their number is at most the smaller of
$|T_2|(|T_2|-1)$ and $2\lfloor(n_0-1)/(4L_w)\rfloor$,
independently of how many rational primes give the same
Frobenius. This is a test in the finite host group.
It uses no density estimate on the primes.
% END H83F2

% BEGIN H83F3
\paragraph{H83F3. Bounding the number of places that can overspend their row.}
Continue with $\psi\ne\mathrm{cc}$. Let $d$ be the size of a
decomposition orbit of $\psi$. The CM host is Galois and
unramified at $p$, so every $u\mid p$ has $d_u=d$.
At most $2M(v)$ such places have a pair of adjacent thick
letters. To count them, assign to each such orbit one starting
embedding of a directed double hit. Different orbits give
different starts. A first digit $s$ specifies one thick real
place and there are just two conjugate embeddings of $E$ above
it, by the CM type description of C82S1. Every start yields the
ordered pair in H83F2. Counting its possible initial embeddings
gives the factor two, also when a cyclic pair crosses the index
zero. A repeated hit within an orbit can only improve this
upper bound.

At all the remaining places take $B_{izu}=b_u$ in E83D1, as
just proved. At each problematic place take the coarse bound of
E83D1. Set $b_{\min}=\min_{u\mid p} b_u$ and
$K_p=(d+1)r(25h_B/p+3)$. The actual signed excess is at most
\[
 \min(m_p,2M(v))\,(K_p-b_{\min})^+
                   -t c_Bm_p-(2n_V+2)\operatorname{vdel}(i).
\]
Indeed a problematic place costs at most $K_p-b_u$ and every
other place has zero difference. If the display is nonpositive,
the original pointwise inequality follows throughout this index.
Otherwise its positive part can be used in E83D2, which counts
only coefficient-field places and good-point multiplicities
after this host-place calculation. All constants here are read
on the same member. In particular no factor for the entire
decomposition degree is inferred from a count of possible
Frobenius elements across different members.
% END H83F3
